\documentclass[12pt,a4paper]{article}
\usepackage[spanish,es-nodecimaldot]{babel}
\usepackage[utf8]{inputenc}
\usepackage[T1]{fontenc}
\usepackage{lmodern}
\usepackage[top=1.1cm,bottom=1.1cm,left=1.6cm,right=1.6cm]{geometry}
\usepackage{amsmath,amssymb}
\usepackage[table]{xcolor}
\usepackage{enumitem}

\setlength{\parindent}{0pt}
\setlength{\parskip}{14pt}
\pagestyle{empty}

\newcommand{\cuerpo}{\fontsize{16.8}{20.5}\selectfont}

\AtBeginDocument{%
  \setlength{\abovedisplayskip}{10pt}%
  \setlength{\belowdisplayskip}{10pt}%
  \setlength{\abovedisplayshortskip}{5pt}%
  \setlength{\belowdisplayshortskip}{5pt}%
  \cuerpo
}

\begin{document}

{\fontsize{18}{22}\selectfont\bfseries Regresion curvilinea\par}

\hspace{16pt}\begin{minipage}{0.94\linewidth}
Ahora que pasa si al graficar el diagrama de dispersion vemos que los datos no siguen una trayectoria recta? queremos ajustar con linea recta pero el modelo no va a ser el mas adecuado ya q no representa la naturaleza del fenomeno. Es por esto q usamos un modelo curvilineo. Para lograr este modelo y que se comporte de manera lineal debemos aplicar una transformacion matematica a los datos originales y asi poder aplicar el metodo de los minimos cuadrados que ya veniamos trabajando.
\end{minipage}

\hspace{16pt}\textbf{Exponencial}

\begin{center}
$y = \alpha\cdot\beta^x$
\quad \textbf{=$>$} \quad
$\log_{10}(y) = \log_{10}(\alpha\cdot\beta^x)$
\quad \textbf{=$>$} \quad
$\log_{10}(y) = \log_{10}(\alpha) + x\cdot\log_{10}(\beta)$

\vspace{8pt}
$y^{*} = \log_{10}(y), \qquad a = \log_{10}(\alpha), \qquad b = \log_{10}(\beta)$

\vspace{8pt}
\textbf{=$>$} \quad $y^{*} = a + b\cdot x$
\hspace{1.6cm}
$\alpha = 10^a$ \hspace{1.0cm} $\beta = 10^b$
\end{center}

\hspace{16pt}\textbf{Potencial}

\begin{center}
$y = \alpha\cdot x^{\beta}$
\quad \textbf{=$>$} \quad
$\log_{10}(y) = \log_{10}(\alpha) + \beta\cdot\log_{10}(x)$

\vspace{8pt}
\textbf{=$>$} \quad $y^{*} = \log_{10}(y), \quad x^{*} = \log_{10}(x), \quad a = \log_{10}(\alpha), \quad B = \beta$

\vspace{8pt}
\textbf{=$>$} \ $y^{*} = a + Bx^{*}$
\hspace{2.2cm}
$\alpha = 10^a$ \hspace{1.0cm} $\beta = 10^b$
\end{center}

\hspace{16pt}\textbf{Reciprocos}

\begin{center}
$y = \dfrac{1}{\alpha+\beta\cdot x}$
\quad \textbf{=$>$} \quad
$\dfrac{1}{y} = \alpha + \beta\cdot x$

\vspace{10pt}
\textbf{=$>$} \ $y^{*} = \dfrac{1}{y}$
\hspace{1.2cm}
$a=\alpha$ \quad $b=\beta$

\vspace{10pt}
\textbf{=$>$} \ $y^{*} = \alpha + \beta x$
\end{center}

\end{document}
